% Stage 8 restored fragment — 2.0.3 Admissibility and Maxwell export sequence (archive lines 917-1113)
\subsection{Admissibility and Maxwell export sequence}
\label{sec:ch203}

A coax-fed pyramidal horn over a ground plane defines the impressed-current and boundary-data tuple technicians declare before Maxwell export is computed \cite{stratton1941}. Maxwell's retarded export to $\Gamma_R$ remains the governing constraint \cite{maxwell1873}.
Subsections~\ref{sec:ch201}--\ref{sec:ch202} introduced the radiation maps and Maxwell
export sequence; the display below invokes \ref{ch:01} $\mathcal{D}_J=1$
\cite{stratton1941}. The present chapter builds on
established in \ref{ch:01}. Rather, it applies those foundations to radiation problems
\cite{maxwell1873}. The Maxwell-closed criteria, conservation laws, and
measurement principles developed previously remain in force; the present chapter introduces
the operator calculus through which admissible sources are connected to observable
radiative fields \cite{harrington2001}. The operators $\mathcal{S}_\omega$,
$\mathcal{R}_\omega$, and $\mathcal{L}_\omega$ are applied on admissible inputs without
reopening the \ref{ch:01} proofs \cite{stratton1941}. Impressed data enter the
solution domain only when the \ref{ch:01} source conditions hold; radiative
reports attach to \(\mathcal{R}_\omega[\Jimp]=\Pi_{\mathrm{rad}}\mathcal{S}_\omega[\Jimp]\),
not to the full solution pair by default \cite{stratton1941}.
The continuity relations below restate, in operator language, the \ref{ch:01}
regime and measurement criteria that qualify when angular-flux and transport reports may
attach to \(\mathcal{R}_\omega[\Jimp]\) rather than to the unreconstructed solution pair
\cite{stratton1941}. Physically, this means that near-field reactive content diagnosed on
\(\mathcal{S}_\omega[\Jimp]\) must not be reported as exported radiative flux unless the
regime screen \(\mathcal{D}_{J,\mathrm{reg}}\) from \ref{ch:01} is satisfied on
the measurement shell \cite{harrington2001}.

\begin{equation}
\label{eq:ch2.continuity-operator-domain}
\mathcal{D}_{\mathrm{ch1}}=1
\ \Longrightarrow\
\Jimp\in\mathcal{J}_\Omega
\ \wedge\
\mathcal{S}_\omega[\Jimp]\in\mathcal{F}_\Omega,
\end{equation}
\begin{equation}
\label{eq:ch2.transport-qualified-radiation}
\mathcal{D}_{J,\mathrm{reg}}=1
\ \Longrightarrow\
\text{angular-flux reports attach to}\ \mathcal{R}_\omega[\Jimp],
\ \text{not to}\ \mathcal{S}_\omega[\Jimp]\ \text{by default},
\end{equation}
% Forward registry: \ref{ch:01} labels imported at Ch.~2 continuity (proofs remain in Ch.~1).
\begin{equation}
\label{eq:ch1.local-regime-decomposition}
\Pi(\theta,\phi;r)
=
\Pi_{\mathrm{a}}(\theta,\phi;r)+\Pi_{\mathrm{re}}(\theta,\phi;r),
\qquad
\Pi_{\mathrm{a}}=\Re\{\mathbf{S}\cdot\hat{\mathbf{n}}\},
\end{equation}
with reactive sector \(\Pi_{\mathrm{re}}\) the remainder on the declared control surface.
\begin{equation}
\label{eq:ch1.reactive-flux-sphere}
Q_{\mathrm{re}}(r)
=
\Im\oint_{S_r}\mathbf{S}\cdot\hat{\mathbf{r}}\,dS,
\qquad
P_{\mathrm{rad}}(r)
=
\Re\oint_{S_r}\mathbf{S}\cdot\hat{\mathbf{r}}\,dS,
\end{equation}
\begin{equation}
\label{eq:ch1.regime-index}
\chi(r)
=
\frac{|Q_{\mathrm{re}}(r)|}{P_{\mathrm{rad}}(r)+\epsilon_P},
\qquad \epsilon_P>0.
\end{equation}
\begin{equation}
\label{eq:ch1.radiative-regime-condition}
P_{\mathrm{rad}}(r)>0
\ \wedge\
\chi(r)\le\chi_{\mathrm{th}},
\end{equation}
\begin{equation}
\label{eq:ch1.crossover-radius}
r_c:\ \chi(r_c)=\chi_{\mathrm{th}},
\end{equation}
\begin{equation}
\label{eq:ch1.regime-qualified-decision}
\mathcal{D}_{J,\mathrm{reg}}(r)=1
\iff
\Big(\chi(r)\le\chi_{\mathrm{th}}\ \wedge\ \mathcal{D}_J(r)=1\Big),
\end{equation}
\begin{equation}
\label{eq:ch1.rotation-power-invariance}
\oint_{S_r}\Re\{\mathbf{S}_R\cdot\hat{\mathbf{r}}\}\,dS
=
\oint_{S_r}\Re\{\mathbf{S}\cdot\hat{\mathbf{r}}\}\,dS,
\end{equation}
\begin{equation}
\label{eq:ch1.scalar-observable-invariance}
\mathcal{I}\big[\mathcal{P}_R\mathcal{F}\big]=\mathcal{I}[\mathcal{F}],
\end{equation}
for scalar observables \(\mathcal{I}\) under passive frame rotation \(\mathcal{P}_R\).
\begin{equation}
\label{eq:ch1.covariance-residual}
\mathcal{R}_{\mathcal{F}}
=
\left\|
\mathcal{R}_\omega\big[\mathcal{T}_{R,\mathbf{a}}\Jimp\big]
-
\mathcal{T}_{R,\mathbf{a}}\,\mathcal{R}_\omega[\Jimp]
\right\|,
\end{equation}
\begin{equation}
\label{eq:ch1.intrinsic-angular-observable}
\mathbf{J}_{\mathrm{int}}
=
\mathbf{J}_{\mathrm{em}}
-
(\mathbf{r}-\mathbf{r}_c)\times\mathbf{P}_{\mathrm{tot}},
\end{equation}
\begin{equation}
\label{eq:ch1.transverse-components}
\E_\perp
=
\E-(\hat{\mathbf{r}}\cdot\E)\hat{\mathbf{r}},
\qquad
\H_\perp
=
\H-(\hat{\mathbf{r}}\cdot\H)\hat{\mathbf{r}},
\end{equation}
\begin{equation}
\label{eq:ch1.stokes-local}
S_0=|E_x|^2+|E_y|^2,\quad
S_1=|E_x|^2-|E_y|^2,\quad
S_2=2\Re(E_x E_y^*),\quad
S_3=2\Im(E_x E_y^*),
\end{equation}
\begin{equation}
\label{eq:ch1.poynting-vorticity}
\Omega_P
=
\hat{\mathbf{n}}\cdot(\nabla\times\mathbf{P}),
\qquad
\mathbf{P}=\Re\{\mathbf{S}\},
\end{equation}
where \(\mathcal{D}_{J,\mathrm{reg}}\) is the regime flag from
Eq.~\eqref{eq:ch1.regime-qualified-decision}.
The laboratory chain composes inherited \ref{ch:01} constraints with radiation
selection and finite observation:
\begin{equation}
\label{eq:ch2.continuity-composition-law}
\mathfrak{A}_{\mathrm{ch2}}^{\mathrm{(report)}}
=
\mathfrak{A}_{\mathrm{meas}}
\circ
\mathfrak{A}_{\mathrm{g}}
\circ
\mathfrak{A}_{\mathrm{cov}}
\circ
\mathfrak{A}_{\mathrm{tr}}
\circ
\Pi_{\mathrm{rad}}
\circ
\mathcal{S}_\omega,
\end{equation}
\begin{equation}
\label{eq:ch2.operational-measurement-chain}
\Lambda_{\mathrm{lab}}
=
\mathcal{O}_{\mathrm{meas}}\circ\mathcal{O}_{\mathrm{far}}\circ\mathcal{R}_\omega,
\qquad
\mathcal{Q}\in\{\text{direct},\ \text{surrogate},\ \text{inference}\},
\end{equation}
% Missing labeled equations restored at natural narrative anchors (math frozen).
\begin{equation}
\label{eq:ch2.covariance-operator-test}
\mathcal{R}_\omega\big[\mathcal{T}_{R,\mathbf{a}}\Jimp\big]
=
\mathcal{T}_{R,\mathbf{a}}\,\mathcal{R}_\omega[\Jimp]
\quad\text{on admissible data,}
\end{equation}
with residual Eq.~\eqref{eq:ch1.covariance-residual} serving as the numerical diagnostic when
discretization breaks Eq.~\eqref{eq:ch2.covariance-operator-test} \cite{stratton1941}.
Function spaces, Green functions, and radiation operators exist because Maxwell field
specification requires them---not as parallel formulations \cite{maxwell1873}. Section~\ref{sec:ch20}
has established the radiation maps, Maxwell export sequence, and \ref{ch:01}
continuity; Section~\ref{sec:ch21} now supplies the Hilbert and Sobolev apparatus through
which those maps acquire precise domains, traces, and adjoints \cite{stratton1941}.
